\subsection{半单环的因子}

在本小节中，我们考虑一个半单环 A。

我们用 \(\mathcal{S}\) 表示单左 A-模的类所成的集（VIII, 第 51 页）；它是有限的（VIII, 第 136 页，命题 1）。对每个 \(\lambda \in \mathcal{S}\)，我们选取一个单 A-模 \(S_{\lambda}\)（类为 \(\lambda\)）。我们用 \(\mathfrak{b}_{\lambda}\) 表示它的零化子，并把它的交换子 \(\mathrm{End}_{\mathrm{A}}(S_{\lambda})\) 的反体记作 \(D_{\lambda}\)。

命题 8. —— a) 对每个 \(\lambda \in \mathcal{S}\)，\(S_{\lambda}\) 是体 \(D_{\lambda}\) 上的一个有限维右向量空间。取商后，映射 \(a \mapsto a_{S_{\lambda}}\) 定义了从 \(A / \mathfrak{b}_{\lambda}\) 到 \(\mathrm{End}_{D_{\lambda}}(S_{\lambda})\) 的环同构。

\begin{enumerate}
\def\labelenumi{\alph{enumi})}
\setcounter{enumi}{1}
\tightlist
\item
  对每个 \(\lambda \in \mathcal{S}\)，环 \(\mathrm{A} / \mathfrak{b}_{\lambda}\) 是单的，且典范同态 \(\psi\)（从 \(\mathrm{A}\) 到 \(\prod_{\lambda \in \mathcal{S}}\mathrm{A} / \mathfrak{b}_{\lambda}\)）是环同构。
\end{enumerate}

环 \(A/b_{\lambda}\) 同构于 \(A_{S_{\lambda}}\)。由 VIII, 第 139 页的引理 1，这个环是单的。所给的从 \(A/b_{\lambda}\) 到 \(\operatorname{End}_{D_{\lambda}}(S_{\lambda})\) 的映射可以等同于从 \(A_{S_{\lambda}}\) 到 \(A_{S_{\lambda}}^{\prime\prime}\) 的映射。由 VIII, 第 139 页的命题 5，它是一个同构。

A-模 \(A_{s}\) 是半单的、忠实的且平衡的。同态 \(\psi\) 可以等同于从 \(A_{s}\) 的双交换子到 \(\prod_{\lambda\in\mathcal{S}}\mathrm{End}_{\mathrm{D}_{\lambda}}(\mathrm{S}_{\lambda})\) 的态射，而后者是同构（VIII, 第 86 页，命题 7, c)）。

单环 \(A/b_{\lambda}\) 称为 A 的 \(\lambda\) 型单因子。

例. —— 设 K 是一个代数闭交换域，A 是 K 上的一个有限次数半单代数。设 \((\mathrm{V}_{i})_{i\in\mathrm{I}}\) 是一个单 A-模族，使得每个单 A-模都同构于唯一的 \(V_{i}\)。则 I 是有限集（VIII, 第 136 页，命题 1），诸向量空间 \(V_{i}\) 在体 K 上都是有限维的，\(V_{i}\) 的交换子等于 \(K\cdot1_{V_{i}}\)（VIII, 第 47 页，定理 1），且映射 \(a\mapsto(a_{V_{i}})_{i\in\mathrm{I}}\) 是从 A 到 \(\prod_{i\in\mathrm{I}}\mathrm{End}_{\mathrm{K}}(\mathrm{V}_{i})\) 的代数同构（命题 8）。

我们已定义（VIII, 第 50 页）极小双边理想为非零双边理想之集（按包含关系排序）中的极小元。换言之，A 的极小双边理想 \(\mathfrak{a}\) 就是 \(s\mathrm{A}_d\) 的单 (A, A)-子双模。同样，我们把 A 的极大双边理想定义为 A 的真双边理想之集中的极大元。A 的极大双边理想 \(\mathfrak{a}\) 就是 \(s\mathrm{A}_d\) 的极大 (A, A)-子双模（VIII, 第 48 页，定义 2）。若环 A 是单的，则理想 0 是 A 的极大双边理想，而 A 是极小双边理想。

对每个 \(\lambda \in \mathcal{S}\)，我们把型为 \(\lambda\) 的同型分量（A-模 \(A_s\) 的）记作 \(\mathfrak{a}_{\lambda}\)。对任一子集 \(\Lambda\)（\(\mathcal{S}\) 的），令 \(\mathfrak{a}_{\Lambda} = \sum_{\lambda \in \Lambda} \mathfrak{a}_{\lambda}\)。

命题 9. —— a) 把集 \(\mathfrak{P}(\mathcal{S})\)（\(\mathcal{S}\) 的子集全体）与集 \(\mathcal{B}_{\mathrm{A}}\)（A 的双边理想全体）按包含关系排序。映射 \(\Lambda \mapsto \mathfrak{a}_{\Lambda}\) 是从 \(\mathfrak{P}(\mathcal{S})\) 到 \(\mathcal{B}_{\mathrm{A}}\) 的有序集同构。

\begin{enumerate}
\def\labelenumi{\alph{enumi})}
\setcounter{enumi}{1}
\item
  A 的极小双边理想是诸理想 \(\mathfrak{a}_{\lambda}\)。
\item
  我们有 \(\mathfrak{b}_{\lambda} = \mathfrak{a}_{\mathcal{S} - \{\lambda\}}\)（对每个 \(\lambda \in \mathcal{S}\)），且诸理想 \(\mathfrak{b}_{\lambda}\) 是 A 的极大双边理想。
\item
  对每个 \(\lambda \in \mathcal{S}\)，从 A 到 \(\mathrm{A} / \mathfrak{b}_{\lambda}\) 的典范映射诱导出从 \(\mathfrak{a}_{\lambda}\) 到 \(\mathrm{A} / \mathfrak{b}_{\lambda}\) 的 A-模同构。
\end{enumerate}

断言 a) 由应用于 A-模 \(\mathbf{A}_s\) 的 VIII, 第 87 页命题 8, d) 得出。由此推出 A 的极小双边理想是诸 \(\mathfrak{a}_{\lambda}\)，而极大双边理想是诸理想 \(\mathfrak{c}_{\lambda} = \mathfrak{a}_{\mathcal{S} - \lambda}\)（对 \(\lambda \in \mathcal{S}\)）。

剩下要证的是 \(b_{\lambda}\) 与 \(c_{\lambda}\) 的相等（对每个 \(\lambda \in S\)）。设 \(\lambda\) 与 \(\mu\) 在 S 中互异。子模 \(a_{\mu}S_{\lambda}\)（A-模 \(S_{\lambda}\) 的）是诸线性映射 \(a \mapsto ax\)（从 \(a_{\mu}\) 到 \(S_{\lambda}\)，\(x \in S_{\lambda}\)）的像的并。因此它为零，且我们有 \(a_{\mu} \subset b_{\lambda}\)。于是我们有 \(c_{\lambda} \subset b_{\lambda}\)，最后 \(c_{\lambda} = b_{\lambda}\)，因为 \(c_{\lambda}\) 是 A 的极大双边理想而 \(b_{\lambda}\) 异于 A。

推论. —— 设 \((\mathrm{A}_{i})_{i\in\mathrm{I}}\) 是单环的一个有限族，f 是从 A 到 \(\prod_{i\in I}A_{i}\) 的一个同构。对每个 \(i\in I\)，存在 S 的唯一元素 \(\varphi(i)\)，使得 \(pr_{i}\circ f\) 的核是 \(\mathfrak{b}_{\varphi(i)}\)。映射 \(\varphi\) 是从 I 到 S 的双射；映射 \(pr_{i}\circ f\) 诱导出一个同构 \(f_{i}\)（从 \(\mathrm{A}/\mathfrak{b}_{\varphi(i)}\) 到 \(A_{i}\)，对每个 \(i\in I\)）。

于是，f 是从 A 到 \(\prod_{\lambda\in\mathscr{S}}A/\mathfrak{b}_{\lambda}\) 的典范同构与从 \(\prod_{\lambda\in\mathscr{S}}A/\mathfrak{b}_{\lambda}\) 到 \(\prod_{i\in I}A_{i}\) 的同构（由诸 \(f_{i}\) 得到）的复合（“半单环分解为单环乘积的唯一性”）。

证明该推论。设 \(i \in I\)；把 \(pr_{i} \circ f\) 的核记作 \(b_{i}^{\prime}\)。由于单环 \(A_{i}\) 同构于 \(A/b_{i}^{\prime}\)，A 的双边理想 \(b_{i}^{\prime}\) 是极大的。于是由命题 8, c)，存在 S 的唯一元素 \(\varphi(i)\)，使得我们有 \(b_{i}^{\prime} = b_{\varphi(i)}\)。取商后，\(pr_{i} \circ f\) 定义了一个同构 \(f_{i}\)（从 \(A/b_{\varphi(i)}\) 到 \(A_{i}\)）。此外，我们有 \(b_{i}^{\prime} + b_{j}^{\prime} = A\)（若 \(i \neq j\)），且 \(\cap_{i \in I} b_{i}^{\prime} = 0\)（参见 I, §8, No.~11, 第 110 页，命题 10）。由此及命题 8 推出 \(\varphi\) 是从 I 到 S 的双射。

命题 10. —— 用 Z 表示 A 的中心。对 \(\lambda \in \mathcal{S}\)，设 \(Z_{\lambda}\) 是体 \(D_{\lambda}\) 的中心。

\begin{enumerate}
\def\labelenumi{\alph{enumi})}
\item
  映射 \(z \mapsto (z_{\mathrm{S}_{\lambda}})_{\lambda \in \mathcal{S}}\) 是从环 \(\mathbf{Z}\) 到乘积 \(\prod_{\lambda \in \mathcal{S}}\mathrm{Z}_{\lambda}\) 的同构。
\item
  把 Z 的理想所成的集 \(I_{Z}\) 与 A 的双边理想所成的集 \(B_{A}\) 按包含关系排序。映射 \(a \mapsto aA\) 是从 \(I_{Z}\) 到 \(B_{A}\) 的有序集同构。逆同构把 A 的双边理想 b 映为 Z 的理想 \(b \cap Z\)。
\end{enumerate}

本命题由应用于 A-模 \(A_{s}\)（其双交换子为 A）的 VIII, 第 87 页命题 8 得出。

推论. —— 设 B 是一个环。下列性质等价：

\begin{enumerate}
\def\labelenumi{(\roman{enumi})}
\item
  环 B 是单的。
\item
  环 B 是半单的，且其中心是一个体。
\item
  环 B 是半单的，且只有一类的单 B-模。
\end{enumerate}

命题 11. —— 设 \(\lambda \in \mathcal{S}\)。A 的同型分量 \(\mathfrak{a}_{\lambda}\) 既是 \(\mathrm{A}_s\) 的 \(\mathrm{S}_{\lambda}\) 型同型分量，又是 \(\mathrm{A}_d\) 的 \(\mathrm{S}_{\lambda}^{*}\) 型同型分量。此外，我们有

\[
[ \mathrm{A} _ {s}: \mathrm{S} _ {\lambda} ] = [ \mathrm{A} _ {d}: \mathrm{S} _ {\lambda} ^ {*} ] = \dim_ {\mathrm{D} _ {\lambda}} \mathrm{S} _ {\lambda}\tag{1}
\]

以及

\[
\operatorname{long} (\mathrm{A}) = \operatorname{long} \left(\mathrm{A} ^ {\mathrm{o}}\right) = \sum_ {\lambda \in \mathscr {S}} \dim_ {\mathrm{D} _ {\lambda}} \mathrm{S} _ {\lambda}.\tag{2}
\]

第一个断言是 VIII, 第 139 页命题 6, c) 当 \(M = A_{s}\) 时的特款。等式 \([A_{s} : S_{\lambda}] = [A_{d} : S_{\lambda}^{*}]\) 由 VIII, 第 140 页的命题 7 得出，因为左 A-模 \(A_{s}\) 的对偶同构于右 A-模 \(A_{d}\)。由命题 8, a) 与 9, c)，映射 \(a \mapsto a_{S_{\lambda}}\) 定义了从 \(a_{\lambda}\) 到 \(\operatorname{End}_{D_{\lambda}}(S_{\lambda})\) 的左 A-模同构。由于按定义 \([A_{s} : S_{\lambda}]\) 是左 A-模 \(a_{\lambda}\) 的长度，关系式 \([A_{s} : S_{\lambda}] = \dim_{D_{\lambda}} S_{\lambda}\) 由 VIII, 第 121 页的引理 2 得出。最后，对 \(\lambda\) 求和，即由 (1) 得到关系式 (2)。

附注。 —— 设 A 是一个半单环，Z 是它的中心。在下列诸集之间存在典范双射：

\begin{enumerate}
\def\labelenumi{\alph{enumi})}
\item
  单左 A-模的类所成的集 \(\mathcal{S}(\mathrm{A})\)
\item
  单右 A-模的类所成的集 \(\mathcal{S}(\mathrm{A}^{\circ})\)
\item
  A 的极小双边理想所成的集
\item
  A 的极大双边理想所成的集
\item
  单 Z-模的类所成的集 \(\mathcal{S}(\mathbf{Z})\)
\item
  Z 的极小理想所成的集
\item
  \(\mathbf{Z}\) 的极大理想所成的集。
\end{enumerate}

于是，对每个元素 \(\lambda\)（属于 \(\mathcal{S}(A)\)），相对应的是类 \(\lambda^{*}\)（单右 A-模 \(S_{\lambda}^{*}\)，\(S_{\lambda}\) 的对偶的）、A 的极小双边理想 \(a_{\lambda}\)（\(A_{s}\) 的 \(\lambda\) 型同型分量）、A 的极大双边理想 \(b_{\lambda}\)（单模 \(S_{\lambda}\) 的零化子）、单 Z-模 \(Z \cap a_{\lambda}\) 的类、Z 的极小理想 \(Z \cap a_{\lambda}\) 以及 Z 的极大理想 \(Z \cap b_{\lambda}\)。

命题 12. —— 设 M 是半单环 A 上的一个模，\(\mathcal{S}_{\mathrm{M}} \subset \mathcal{S}\) 是 M 的支集。则 M 的零化子 Ann(M) 是双边理想 \(\sum_{\lambda \in \mathcal{S} - \mathcal{S}_{\mathrm{M}}} \mathfrak{a}_{\lambda}\)，而 M 的迹理想 \(\tau(\mathrm{M})\) 是双边理想 \(\sum_{\lambda \in \mathcal{S}_{\mathrm{M}}} \mathfrak{a}_{\lambda}\)。特别地，A 是 Ann(M) 与 \(\tau(\mathrm{M})\) 的直和。

按定义（VIII, 第 84 页），\(\mathcal{S}_{\mathrm{M}}\) 由 M 的单子模的类组成。由于模 M 是半单的，M 的零化子是零化子 \(\mathfrak{b}_{\lambda}\)（对应于类为 \(\lambda\) 的模，\(\lambda\) 取遍 \(\mathcal{S}_{\mathrm{M}}\)）的交。而我们有 \(\mathfrak{b}_{\lambda} = \sum_{\mu \neq \lambda} \mathfrak{a}_{\mu}\)（对每个 \(\lambda \in \mathcal{S}\)）（命题 9）。由于 A 是族 \((\mathfrak{a}_{\lambda})_{\lambda \in \mathcal{S}}\) 的直和，M 的零化子确实等于 \(\sum_{\lambda \in \mathcal{S} - \mathcal{S}_{\mathrm{M}}} \mathfrak{a}_{\lambda}\)。

按定义（VIII, 第 80 页），迹理想 \(\tau(\mathrm{M})\) 是由从 M 到 \(\mathrm{A}_s\) 的 A-线性映射的像生成的 \(\mathrm{A}_s\) 的 A-子模。由于 M 是半单的，这等价于说 \(\tau(\mathrm{M})\) 由 \(\mathrm{A}_s\) 的类属于 \(\mathcal{S}_{\mathrm{M}}\) 的单子模生成。于是我们有 \(\tau(\mathrm{M}) = \sum_{\lambda \in \mathcal{S}_{\mathrm{M}}} \mathfrak{a}_{\lambda}\)。

推论. —— 设 M 是半单环 A 上的一个模。下列性质等价：

\begin{enumerate}
\def\labelenumi{(\roman{enumi})}
\item
  A-模 M 是忠实的。
\item
  M 的支集等于 S。
\item
  A-模 M 是生成模。
\end{enumerate}

事实上，说 M 是忠实的，即是说它的零化子约化为 0；而 M 是生成模当且仅当它的迹等于 A（VIII, 第 80 页，定理 1）。

